% Authors' Introduction, supplied 2026-08-26. Paragraphs 1-3 are kept verbatim apart from
% added citations and the corrected DNP3 framing/CRC explanation. The contribution
% list was rewritten to verb-led form without a firstness claim. See
% pipeline/reports/LIN_TEXT_CHANGELOG.md.
\section{Introduction}
\label{sec:intro}

Device fingerprinting has been an essential step in cyber reconnaissance, allowing adversaries
to reveal unique features of target networks and design effective, stealthy attack strategies.
These techniques are becoming increasingly critical in industrial control systems (ICSs) such as
power grids, where adversaries often use IP-based control networks and computing devices within
as entry points to inflict physical damage. Consequently, fingerprinting shifts the focus from
visited websites and user biometric behavior to device models and types of control operations
that are critical to ICS attacks. In the 2015 attack that disrupted Ukrainian power grids and
the Stuxnet attack that disrupted Iranian nuclear power facilities, it is widely believed that
adversaries stay in their systems for at least 6 months to perform cyber
reconnaissance~\cite{leeAnalysisCyberAttack2016, langnerStuxnetDissectingCyberwarfare2011}.

To disrupt device fingerprinting, many studies present network traffic
obfuscation~\cite{wrightTrafficMorphingEfficient2009, dyerPeekaBooStillSee2012}. Because device
fingerprinting targeting general computing environments generally relies on network-level
features such as packet size and/or inter-packet latency observed from communication
patterns~\cite{sirinamDeepFingerprintingUndermining2018}, traffic obfuscation often focuses on
(i) padding and splitting network packets, which hide or change the distribution of network
packet sizes~\cite{wrightTrafficMorphingEfficient2009}, and (ii) delaying network packets and
adding dummy ones, which disrupt the inter-packet timing pattern~\cite{dyerPeekaBooStillSee2012,
caiSystematicApproachDeveloping2014}. Since manipulating communication networks can introduce
runtime overhead, recent works have begun to offload traffic obfuscation onto programmable
network switches, which change communication patterns at much higher line rates than
CPUs~\cite{meierDittoWANTraffic2022, wangMinosLightweightDynamic2025,
xieSecuritasDefendingTraffic2026}.

Unfortunately, these methods do not directly transfer device fingerprinting in ICS/SCADA
environments for two main reasons. First, the leaked features are different. General web traffic
obfuscation usually targets the flow level sizes and timing patterns~\cite{meierDittoWANTraffic2022}
whereas ICS device fingerprinting is carried by the response timing of the protocol exchanges,
which stays visible even when the payload is encrypted~\cite{formbyWhosControlYour2016}.
Secondly, adding bytes or cover packets is constrained in this setting in a way it is not on the
open Internet. A DNP3 link-layer frame carries a declared length and a cyclic redundancy check
(CRC) over each block of the frame~\cite{ieeeIEEEStandardElectric2012}, so bytes inserted at an
arbitrary offset change the framing the receiver parses. Making such an insertion valid requires
protocol-aware reconstruction of the frame and recomputation of its link-layer CRC blocks, which
is not what an opaque pad does. The outstation that would otherwise emit cover traffic is a
protective relay whose firmware cannot be modified, so a legacy relay cannot simply be made to
generate that traffic either. Consequently, a defense in this setting must hide the timing that
leaks even after encryption while changing no bytes.

In this paper, we present an in-network mechanism that normalizes an outstation's
fingerprintable features using the data plane of the Tofino programmable switch at the
outstation edge, without changing the endpoints or modifying the bytes of the DNP3 traffic. This
defense holds per transaction responses and releases them at configurable offsets so the
master-visible timing then becomes a policy value. The two timing features that Formby et
al.~\cite{formbyWhosControlYour2016} use for fingerprinting, the cross-layer response time of a
poll and the master-visible timing of a control command, are the two case studies through which
we show the framework working; they are not two separate systems. We implement the framework as
a P4 program on a single switch, evaluate it on a cyberphysical testbed, and make the following
contributions:
\begin{itemize}
  \item \textbf{Design an in-network timing-obfuscation framework} that replaces selected
    within-transaction timing features with configurable release-policy values, without
    modifying either endpoint and without changing the protected DNP3 bytes.
  \item \textbf{Implement the framework on an Intel Tofino-1} as a P4 program with a separate
    acknowledgment-anchored read lane and request-anchored control lane, bounded queues, and
    fail-open forwarding when a response arrives after its scheduled release.
  \item \textbf{Evaluate the framework against a physical SEL-751A relay} over 22 grouped
    collection runs, and report policy programmability, timing-feature suppression, the effect
    on a fixed transaction-class attacker, the leakage that remains against an adaptive
    attacker, and the latency the mechanism costs.
\end{itemize}
